\section{Further research questions}\label{tsc:further}
\subsection{Joint triple collisions and normalization conversion}
For two approaching seams, the separated three-factor formula develops a
local singularity. When all three approach one another, two relative
scales are available. Determine the full two-parameter local expansion,
including the delta derivatives and finite constants needed to compare
separated analytic continuation with a chosen same-point Hadamard product.
A first concrete target is the case $(m_1,m_2,m_3)=(0,0,0)$ along
$a_2-a_1=\lambda\varepsilon$, $a_3-a_1=\varepsilon$ with fixed
$\lambda\notin\{0,1\}$. The dependence on $\lambda$ must be computed,
not erased by an unproved path-independent finite-part prescription.

\subsection{Transverse Tornheim jets and genuine arithmetic reductions}
Proposition~\ref{tsc:negative-resonance} gives finite reductions on integral
spectral hyperplanes, but the Stieltjes formula also uses derivatives
transverse to them. Find additional functional identities that reduce
specific transverse jets at rational colors. The question should be asked
first at a fixed conductor and a fixed total spectral degree. A negative
result must distinguish nonmembership in an explicitly generated formal
relation space from independence of the resulting numerical periods.

\subsection{Four unreflected factors and balanced frequency sectors}
Four nonzero frequencies summing to zero have sectors of types $3+1$ and
$2+2$. The latter are not a direct copy of the Tornheim double sum. Develop
an explicit, jointly regularized coordinate family for these sectors and
a counterpart of Theorem~\ref{tsc:mellin-all}. A useful first target is
four digamma factors at four distinct rational shifts, with every endpoint
constant and contact correction specified. The reflected finite-depth
formula already proved here should serve as a stringent specialization
check, not as evidence that the generic family has depth one.

\subsection{Rational traces that create collisions}
Integer dilation and conductor operations can turn previously distinct
shifts into coincident shifts. Determine the exact trace identities for
the three-factor kernel across that boundary. The existing bilinear
conductor laws and the separated lowering law are useful inputs, but they
do not by themselves determine the new collision counterterms.

\subsection{Certified polylogarithmic jet evaluation}
The present stable series evaluation avoids a known source of numerical
fragility: differentiating a polylogarithm routine at an infinitesimal
order increment. Replace its floating-point truncations with explicit
complex interval enclosures, including a uniform Jonqui\`ere-series
remainder for the unit-color region. Combine that with the subtracted
Mellin integral to obtain rigorous numerical certificates for generic
three-digamma values. This is an algorithmic refinement of proved
identities, not a substitute for their analytic proof.

\subsection{A precise bridge to the remaining fixed-weight candidates}
The repository's $S_6$ and $S_8$ questions concern particular polylogarithmic
period reductions. The present triple calculus does not automatically
settle them. Search for an exact transformation that expresses a candidate's
remaining obstruction in one of the rational-color transverse jets above,
with all boundary terms retained. Such a bridge, together with a proved
jet reduction, would be a meaningful route to a certificate. Numerical
agreement or a shared use of polylogarithms is not such a bridge.

\subsection{Formalization with analytic hypotheses exposed}
A useful formal development would separate the polynomial-growth Fourier
construction, the local meromorphic extension of $x_+^{u-r-1}$, the
separated product lemma, and the finite coefficient algebra. The latter
already has executable exact tests, but not proof-assistant certificates.
The first analytic formal target can be
\eqref{tsc:single-contact}; the trilinear master theorem can then be proved
conditional on the ordinary absolutely convergent Fourier identity before
formalizing the entire continuation.

\section*{Conclusion}
\addcontentsline{toc}{section}{Conclusion}
Three separated Stieltjes singularities have an exact, all-index colored
Tornheim description, and their Laurent coefficients have a convergent
polylogarithmic Mellin representation with explicit endpoint constants.
Mean-zero integration and contact-corrected differentiation extend that
same kernel to Gamma and polygamma products. Reflection selects an
arbitrary-depth subclass with finite single-Stieltjes, and at rational
index-zero shifts finite Gamma, reductions. The resolved coordinate-class
question is therefore a concrete advance within the inspected research
programme, while collisions, transverse arithmetic reductions, and the
remaining fixed-weight conjectures retain clearly stated open targets.
